\section{Introduction}\label{sec:introduction}

Urban air pollution remains a persistent public-health and environmental problem, especially in dense metropolitan areas where traffic emissions, atmospheric conditions, and land-use patterns interact over time \cite{Dominski2021}. Low-emission zones (LEZs), pedestrianization, and related mobility restrictions have therefore become common policy tools for reducing exposure to traffic-related pollutants. Their evaluation is methodologically demanding because air quality changes are shaped not only by regulation but also by weather, seasonality, traffic redistribution, and local urban form.

This study approaches that problem through predictive counterfactuals. The core idea is to estimate forecasting models using only the pre-policy regime and then compare their expected pollution levels with the realized observations after the policy boundary. Such a design is attractive because it can exploit rich multivariate time series, but it is only persuasive when the substantive conclusion survives reasonable choices about model class, temporal context, input information, and placebo boundaries. A positive result observed under a single architecture or a single preprocessing recipe is difficult to defend as policy evidence.

The article therefore asks a narrower but more defensible question: do post-policy pollution levels in the selected Madrid districts remain systematically below the no-change forecast under a compact but robust forecasting design? The emphasis is on stability rather than on rhetorical force. The goal is not to claim definitive causal identification, but to test whether the post-policy deviation persists when the counterfactual is confronted with stronger baselines, placebo cuts, and lightweight sensitivity checks.

The manuscript makes four concrete contributions. First, it compares Ridge, LightGBM, and LSTM against strong naive benchmarks. Second, it studies temporal context through 6- and 24-hour main windows plus a supplementary 12-hour rolling-origin sensitivity. Third, it complements PIP with mean and median post-policy residuals, bootstrap intervals, and late-2021 placebo boundaries. Fourth, it uses feature ablations, a direct-evaluation sensitivity, and explicit local framing for the three-district panel rather than suggesting city-wide generalization.
